% Options for packages loaded elsewhere
\PassOptionsToPackage{unicode}{hyperref}
\PassOptionsToPackage{hyphens}{url}
%
\documentclass[
  11pt,
]{article}
\usepackage{amsmath,amssymb}
\usepackage{iftex}
\ifPDFTeX
  \usepackage[T1]{fontenc}
  \usepackage[utf8]{inputenc}
  \usepackage{textcomp} % provide euro and other symbols
\else % if luatex or xetex
  \usepackage{unicode-math} % this also loads fontspec
  \defaultfontfeatures{Scale=MatchLowercase}
  \defaultfontfeatures[\rmfamily]{Ligatures=TeX,Scale=1}
\fi
\usepackage{lmodern}
\ifPDFTeX\else
  % xetex/luatex font selection
\fi
% Use upquote if available, for straight quotes in verbatim environments
\IfFileExists{upquote.sty}{\usepackage{upquote}}{}
\IfFileExists{microtype.sty}{% use microtype if available
  \usepackage[]{microtype}
  \UseMicrotypeSet[protrusion]{basicmath} % disable protrusion for tt fonts
}{}
\makeatletter
\@ifundefined{KOMAClassName}{% if non-KOMA class
  \IfFileExists{parskip.sty}{%
    \usepackage{parskip}
  }{% else
    \setlength{\parindent}{0pt}
    \setlength{\parskip}{6pt plus 2pt minus 1pt}}
}{% if KOMA class
  \KOMAoptions{parskip=half}}
\makeatother
\usepackage{xcolor}
\usepackage[margin=1in]{geometry}
\setlength{\emergencystretch}{3em} % prevent overfull lines
\providecommand{\tightlist}{%
  \setlength{\itemsep}{0pt}\setlength{\parskip}{0pt}}
\setcounter{secnumdepth}{-\maxdimen} % remove section numbering
\ifLuaTeX
  \usepackage{selnolig}  % disable illegal ligatures
\fi
\usepackage{bookmark}
\IfFileExists{xurl.sty}{\usepackage{xurl}}{} % add URL line breaks if available
\urlstyle{same}
\hypersetup{
  pdftitle={Unnormalized relative base energy at a fibred zero-volume class},
  pdfauthor={Xean research manuscript},
  hidelinks,
  pdfcreator={LaTeX via pandoc}}

\title{Unnormalized relative base energy at a fibred zero-volume class}
\usepackage{etoolbox}
\makeatletter
\providecommand{\subtitle}[1]{% add subtitle to \maketitle
  \apptocmd{\@title}{\par {\large #1 \par}}{}{}
}
\makeatother
\subtitle{Partial result for P15.1; bounded potentials and arbitrary
calibrated Kähler directions}
\author{Xean research manuscript}
\date{4 October 2026}

\begin{document}
\maketitle

\subsection{Status and exact scope}\label{status-and-exact-scope}

This is a \textbf{partial} result for P15.1. It proves the
supplied-fibration construction below, including bounded potentials,
perturbative potentials that need not descend, uniform convergence with
no rate relative to the perturbation parameter, and critical or
higher-dimensional fibres. It is not a solution for an arbitrary
pseudoeffective zero-volume class. No result about P15.3 is claimed. The
measure statements below are supporting estimates for the energy and its
first variation, not a claim to solve the separate general P15.2 target.
No literature-priority claim is made.

The proof uses the precisely stated standard foundations and the new
applications proved below. The unverified operator proposal and earlier
comparison papers supply no mathematical premise. Original research
inputs and the complete frozen argument remain in the supporting packet.
An independent Xean review accepted exactly the bounded supplied-fibration
scope, including the optional continuous-descent assertion at its separate
connected-fibre hypothesis.

Throughout, \[dd^c=\frac{i}{2\pi}\partial\bar\partial,
\qquad d^c=\frac{i}{4\pi}(\bar\partial-\partial).\] In particular, for a
real smooth function \(w\),
\(dw\wedge d^cw=(i/2\pi)\partial w\wedge\bar\partial w\geq0\).

Let \(\pi:X\longrightarrow B\) be a supplied surjective holomorphic map
of connected smooth compact Kähler manifolds, with
\[\dim_{\mathbb{C}}X=n,\quad\dim_{\mathbb{C}}B=\nu,
\quad 1\leq\nu<n,\quad m=n-\nu.\] There is no assumption of a product,
flatness, equidimensionality, or smooth fibres. Connected fibres are
used only for the optional continuous-descent assertion proved below.
Fix a Kähler form \(\beta\) on \(B\), and put
\[\theta=\pi^*\beta,\quad \alpha=[\theta],\quad
V_B=\int_B\beta^\nu>0,\quad K=\binom n\nu.\] A bounded \(q\)-psh
function means an everywhere finite, globally bounded, upper
semicontinuous quasi-psh representative \(v\) satisfying
\(q+dd^cv\geq0\). Norms are pointwise suprema, not essential suprema
with respect to a measure. All products of bounded-potential currents in
this paper are ordinary Bedford--Taylor (BT) products.

For a smooth Kähler form \(q\) on a compact \(d\)-dimensional manifold
\(Y\), set \[T_v=q+dd^cv,\qquad
E_q(v)=\frac1{d+1}\sum_{j=0}^d\int_Yv\,T_v^j q^{d-j},\] with reference
potential \(0\). Its relative form is \[R_q(v,w)=\frac1{d+1}\sum_{j=0}^d
\int_Y(v-w)T_v^jT_w^{d-j}.\] The same expressions on \(B\) use
\(q=\beta\), \(d=\nu\), and \(C_\varphi=\beta+dd^c\varphi\).

\subsection{Precise foundations and source
audit}\label{precise-foundations-and-source-audit}

The only pluripotential foundations imported for the main result are the
following bounded statements.

\textbf{BT foundation.} On a complex manifold, locally bounded psh
potentials \(p_1,\ldots,p_r\) define a closed positive product
\(dd^cp_1\wedge\cdots\wedge dd^cp_r\) by the recursion \(dd^c(p_1 S)\)
for the remaining closed product \(S\). These products are symmetric and
depend only on the currents. If, on coordinate domains, each \(p_{i,k}\)
decreases to a locally bounded psh \(p_i\), the mixed products converge
weakly to the corresponding mixed product. One may use smooth local psh
decreasing regularizations on smaller domains. These assertions concern
locally bounded potentials and local regularizations; they do not assert
global smooth regularization in a fixed Kähler class. Citation: E.
Bedford and B. A. Taylor, \emph{A new capacity for plurisubharmonic
functions}, Acta Math. 149 (1982), 1--40, Section 2; see also \emph{Fine
topology, Shilov boundary, and \((dd^c)^n\)}, J. Funct. Anal. 72 (1987),
225--251. The bounded construction is expressly stated in BEGZ,
arXiv:0812.3674v2, LF175--188, PDF pp.~4--5; the local decreasing
approximation/continuity application also appears in its proof of
Theorem 1.14, LF668--679. We use this standard published bounded
theorem, not an unstated theorem for unbounded products.

Symmetry and multilinearity needed here can also be checked directly
from this foundation: choose simultaneous smooth decreasing local
regularizations; the smooth identities hold at every level; each term
has a BT limit. Differences of positive bounded-potential currents are
treated by finite expansion into positive products. Adding local smooth
potentials for smooth closed forms gives the same construction and
identifies products on chart overlaps.

\textbf{Bounded integration by parts.} On a compact Kähler manifold
\(Y\), let \(\Theta\) be a closed positive current of bidimension
\((1,1)\). Let \(a_1,a_2,b_1,b_2\) be globally bounded quasi-psh
functions, and put \(a=a_1-a_2\), \(b=b_1-b_2\). In the
bounded-potential products used here, \[\int_Y a\,dd^cb\wedge\Theta
=\int_Y b\,dd^ca\wedge\Theta
=-\int_Y db\wedge d^ca\wedge\Theta. \tag{IBP}\] The signed measures are
finite; the gradient products are defined by BT polarization, and
\(da\wedge d^ca\wedge\Theta\) is a positive finite measure. Citation: S.
Boucksom, P. Eyssidieux, V. Guedj, A. Zeriahi, \emph{Monge--Ampère
equations in big cohomology classes}, Acta Math. 205 (2010), 199--262,
Theorem 1.14 and Lemma 1.15; the primary copy is arXiv:0812.3674v2,
LF626--762, PDF pp.~14--16. The statement just used is its special case
with the closed complete pluripolar set \(A=\varnothing\). At every
application below, \(\Theta\) is a mixed BT product of bounded-potential
positive currents and smooth semipositive forms, is closed and positive
by the BT foundation, and has exactly bidimension \((1,1)\). All
functions are globally bounded quasi-psh or explicitly their
differences. In particular, no boundary around a critical fibre or
singularity locus is discarded.

For completeness, the positivity convention for the gradient square is
given locally by \[dd^c p^2=2\,dp\wedge d^cp+2p\,dd^cp\] for a
nonnegative locally bounded psh function \(p\) and a closed positive
\(\Theta\), followed by polarization. The finite-integral and global
integration-by-parts conclusion is precisely the bounded special case of
the cited theorem; it is not inferred by taking weak limits against
arbitrary discontinuous test functions.

\textbf{Nef-volume convention.} For nef real \((1,1)\) classes on a
compact Kähler manifold, their positive intersection product is the
ordinary cup product. Their volume is the top positive product,
including the boundary limit from Kähler perturbations. We use the nef
numerical dimension
\[\operatorname{nd}(\alpha)=\max\{k:\int_X\alpha^k[\kappa]^{n-k}>0\}\]
for a Kähler form \(\kappa\). Citation: BEGZ, Definition 1.17 and
subsequent discussion, arXiv:0812.3674v2, LF810--840, PDF pp.~17--18
(and the usual nef numerical-dimension convention). The classes used
below are nef because they have a smooth semipositive representative,
and their positive perturbations have smooth Kähler representatives with
minimal singularities. Thus this convention applies. The actual
intersections and perturbation volumes are computed in the proof.

\textbf{Ordinary differential foundations.} We use Stokes' theorem for
currents, and Sard's theorem: the critical values of a smooth map
between finite-dimensional smooth manifolds, with smoothness exceeding
the dimension threshold, have measure zero. Citation for Sard: A. Sard,
\emph{The measure of the critical values of differentiable maps}, Bull.
Amer. Math. Soc. 48 (1942), 883--890. Here the holomorphic map is
smooth, and surjectivity onto the positive-dimensional manifold \(B\)
ensures a nonempty regular-value set and a nonempty full-rank open set.
Fibre integration at a regular value is ordinary integration in local
submersion coordinates and a partition of unity, not a singular-fibre
theorem.

The optional continuous-descent paragraph imports one additional
standard foundation, stated with its exact needed scope there. It is not
needed for the main construction on supplied pullbacks.

The primary text was extracted by LF bytes, not \texttt{splitlines};
internal formfeeds were retained. All six supplied excerpt hashes match
exactly. The complete text has SHA256
\texttt{9d52c95247fd7f69194c84df01ef54b4cd98b6e4e35b4f3851f90f04c98a71b0};
the primary PDF has SHA256
\texttt{43097be8459edfd948929fade34367fa7bbd0fe691d9cba1a4f27a11181214cd}.
The supplied LF175--188, LF189--242, LF249--322, LF626--762, LF810--840,
and LF1120--1222 bodies were read and preserved. Their full provenance
is in \texttt{source-provenance.json}.

No proof is borrowed from Proposition 1.4's sentence, ``The proof of the
last point is similarly easy but tedious and will be left to the
reader.'' Multilinearity in our bounded setting follows from the local
regularization argument above. The printed/transcribed proof of Theorem
1.14 contains the repeated-target phrase
``\(v d^c u_k\wedge\Theta\to v d^c u_k\wedge\Theta\)'' and a lowercase
\(\theta\) in a Cauchy--Schwarz line. These are preserved in the
excerpt; they are not silently represented as correct formulas. The
weighted-energy derivative display in Proposition 2.8 suppresses the
definition's \(1/(n+1)\) factor; its sign conclusion is unaffected, but
our exact energy identities are derived independently below. The
transcribed Lemma 1.5 does not explicitly require increasing or eventual
exhaustion of its Borel sets; its displayed dominated-convergence proof
requires such a condition. We do not use that lemma, any unbounded
truncation limit, or the global regularization assertion inside Lemma
1.15's proof. The indispensable imported assertion is the standard
bounded IBP theorem as explicitly stated above.

\subsection{The theorem}\label{the-theorem}

Let \(\omega_\epsilon\) be Kähler forms on \(X\) converging smoothly, as
\(\epsilon\downarrow0\), to a Kähler form \(\omega_0\). A fixed
arbitrary Kähler form is included. Define
\[c_\epsilon=\pi_*(\omega_\epsilon^m),\quad
q_\epsilon=\theta+\epsilon\omega_\epsilon,\quad
s_\epsilon=Kc_\epsilon\epsilon^m.\] Here the degree-zero pushforward is
a strictly positive constant, as proved below. For a fixed form write
\(c_\omega\) instead. If \(\varphi\) is any bounded \(\beta\)-psh
function, set \(U=\varphi\circ\pi\). Let \(u_\epsilon\) be any globally
bounded \(q_\epsilon\)-psh functions with
\[\delta_\epsilon=\|u_\epsilon-U\|_\infty\longrightarrow0.\] No
descending assumption on \(u_\epsilon\), no monotonicity in
\(\epsilon\), and no bound on \(\delta_\epsilon/\epsilon\) are imposed.

Then
\[\lim_{\epsilon\downarrow0}\frac{E_{q_\epsilon}(u_\epsilon)}{s_\epsilon}
=E_\beta(\varphi)
=\frac1{\nu+1}\sum_{a=0}^\nu\int_B\varphi\,C_\varphi^a\beta^{\nu-a}. \tag{1}\]
More generally, if \(v_\epsilon\) uniformly approximates
\(\psi\circ\pi\) for bounded \(\beta\)-psh \(\psi\), and is bounded and
\(q_\epsilon\)-psh, then
\[\lim_{\epsilon\downarrow0}\frac{R_{q_\epsilon}(u_\epsilon,v_\epsilon)}{s_\epsilon}
=R_\beta(\varphi,\psi). \tag{2}\] One may instead calibrate by
\(Kc_{\omega_0}\epsilon^m\), since \(c_\epsilon\to c_{\omega_0}>0\).

The supporting first-variation measures satisfy
\[\frac{\pi_*\big((q_\epsilon+dd^cu_\epsilon)^n\big)}{s_\epsilon}
\ \rightharpoonup\ C_\varphi^\nu. \tag{3}\] The limiting measure has raw
mass \(V_B\). The energy has reference \(0\), value \(E_\beta(0)=0\),
and constant-shift law
\[E_\beta(\varphi+b)=E_\beta(\varphi)+bV_B. \tag{4}\] It is monotone,
concave, and \(V_B\)-Lipschitz for the uniform norm on bounded
\(\beta\)-psh potentials. Its relative form is a cocycle and has the
exact formula already displayed for \(R_\beta\). Its variational scopes
are bounded affine segments and continuously differentiable paths in the
uniform norm of continuous \(\beta\)-psh functions, proved below. No
differentiability for arbitrary bounded time-dependent paths is
asserted.

The conclusions are the same for every Kähler representative and every
Kähler class of the auxiliary direction, after its own specified \(c\)
calibration. The constant \(c\) itself is independent of the
representative of a fixed class but generally changes with the class. No
arbitrary degenerate auxiliary path is included.

The retained data are the supplied carrier \(B\) of dimension \(\nu\),
the map \(\pi\), the base Kähler class and chosen reference current
\(\beta\), the zero reference potential (including its constant), and
the bounded potential/singularity data. The raw reference mass is
\(V_B\), not one. The quantity \(s_\epsilon\) is a positive transverse
calibration; it is not division by \(\operatorname{vol}(\alpha)=0\). The
full leading class-volume coefficient is \(Kc_{\omega}V_B\), and
dividing by that coefficient would lose the stated raw base scale.

\subsection{Pushforward, bounded projection, and numerical
data}\label{pushforward-bounded-projection-and-numerical-data}

First fix any Kähler form \(\omega\). Pushforward of currents is defined
by \[\langle\pi_*S,\eta\rangle=\langle S,\pi^*\eta\rangle\] for
compactly supported smooth test forms on \(B\). Properness follows from
compactness. For \(S=\omega^m\) this is a positive current of bidegree
\((0,0)\), since pullback of positive top-degree test forms and wedging
with \(\omega^m\) is positive. Stokes gives \(d\pi_*S=\pi_*dS=0\). A
degree-zero distribution with all distributional derivatives zero is
constant on a connected manifold: on a ball convolve with a smooth
approximate identity, obtaining functions with zero gradient, and pass
to the distributional limit; constants agree on overlapping balls. Thus
\(\pi_*\omega^m=c_\omega\) on \textbf{all} of \(B\). This also rules out
an additional distribution supported on critical values.

There is a full-rank point of \(\pi\). Otherwise every value would be
critical, contrary to Sard's theorem and surjectivity. Consequently
\(\theta^\nu\omega^m\) is strictly positive on a nonempty open set and
nonnegative everywhere. Pairing the constant current with \(\beta^\nu\)
proves \[c_\omega=\frac{\int_X\theta^\nu\omega^m}{V_B}>0. \tag{5}\] At
any regular value \(b\), local submersion fibre integration identifies
it with \[c_\omega=\int_{\pi^{-1}(b)}\omega^m,\] including the sum over
all components if the regular fibre is disconnected. No assertion that
each singular fibre's naive integral is this number is needed. Formula
(5) gives its cohomological meaning and proves fixed-class
representative independence by Stokes. Scaling the auxiliary form gives
\(c_{t\omega}=t^m c_\omega\), so class independence of \(c\) would be
false.

Here is a bounded projection identity that does not assume smooth fibres
or a preexisting pullback of arbitrary singular currents. On a base
coordinate neighbourhood choose locally bounded psh potentials
\(p_1,\ldots,p_k\) for positive bounded-potential currents
\(C_1,\ldots,C_k\). Define their pullback product by the BT products of
\(p_i\circ\pi\). Holomorphic composition is psh, and local boundedness
is retained. Local decreasing smooth psh approximations \(p_{i,l}\) on
smaller base domains pull back to decreasing smooth psh approximations
on their preimages. BT continuity shows that the definition agrees with
the usual smooth pullback limit and is independent of potentials. Smooth
products of more than \(\nu\) base factors vanish before taking this
limit; hence the same is true for their bounded pullback products, even
at critical points.

For \(k\leq\nu\), the exact projection identity is
\[\pi_*\left(\pi^*C_1\wedge\cdots\wedge\pi^*C_k\wedge\omega^m\right)
=c_\omega C_1\wedge\cdots\wedge C_k. \tag{6}\] To prove it without any
discontinuous-weight weak-limit inference, put \(P_0=\omega^m\),
\(P_r=dd^c((p_r\circ\pi)P_{r-1})\). These are closed positive BT
currents. For any order-zero current \(S\) and bounded Borel function
\(g\) on the base, \[\pi_*((g\circ\pi)S)=g\,\pi_*S\] whenever the
indicated products have order zero: test in local coordinates, integrate
the coefficient measures against \(g\circ\pi\), and use the definition
of their pushforward. In our use the currents are positive, so
coefficient measures are dominated by the trace measure and these
bounded products are defined. Proper pushforward also commutes with
\(dd^c\). Inductively, \[\pi_*P_r=dd^c(p_r\pi_*P_{r-1})
=c_\omega dd^cp_r\wedge\cdots\wedge dd^cp_1.\] The recursion uses
exactly the bounded BT definition. This proves (6) locally on every base
chart; the identities patch. It is an equality on all of \(B\), not
merely off the critical values. In particular, for any bounded Borel
\(g\) and \(a+b=\nu\),
\[\int_X(g\circ\pi)(\theta+dd^cU)^a\theta^b\omega^m
=c_\omega\int_Bg\,C_\varphi^a\beta^b. \tag{7}\] Take \(g=\varphi\) for
the energy. This proves the weighted projection required for possibly
discontinuous bounded \(\varphi\), without testing a weak limit against
\(\varphi\). Smooth factors \(\beta\) are included using their smooth
local psh potentials. Identity (6) is claimed only for bounded local
potentials and the specified smooth transverse form; no arbitrary
singular-current pullback/projection theorem is assumed.

We will repeatedly use the following cohomological mass identity. If
\(Q_i=q_i+dd^cw_i\geq0\) have bounded global potentials and smooth
closed backgrounds, their mixed BT product has cohomology class
\([q_1]\cdots[q_r]\). Indeed, replacing one factor at a time changes the
product by \[dd^c(w_i S),\] where \(S\) is the remaining closed product.
This identity is local BT recursion and patches globally; \(w_iS\) has
order zero because \(w_i\) is bounded. Pairing with any smooth closed
complementary form gives zero by Stokes. The argument also proves the
top-degree mixed mass identities, including those involving \(\theta\),
and does not require strict positivity of \(\theta\). All positive
top-degree products below are finite measures on compact \(X\) or \(B\).

Since \(\theta\geq0\) is smooth and \(\theta+\epsilon\omega>0\),
\(\alpha\) is nef and pseudoeffective. One has \(\theta^{\nu+1}=0\) and,
by (5), \(\int_X\theta^\nu\omega^m=c_\omega V_B>0\). The stated nef
convention therefore gives \(\operatorname{nd}(\alpha)=\nu\) and
\[\operatorname{vol}(\alpha)=\int_X\theta^n=0.\] For every Kähler
direction, its perturbative volume is exactly
\[\int_X(\theta+\epsilon\omega)^n
=\sum_{k=0}^\nu\binom nk\epsilon^{n-k}\int_X\theta^k\omega^{n-k}
=Kc_\omega V_B\epsilon^m+O(\epsilon^{m+1}). \tag{8}\] This also directly
verifies the boundary volume limit. Smooth convergence of
\(\omega_\epsilon\) makes these finitely many coefficient integrals
uniformly bounded, and (5) shows \(c_\epsilon\to c_{\omega_0}>0\).

\subsection{Bounded energy identities, with their actual variational
scopes}\label{bounded-energy-identities-with-their-actual-variational-scopes}

Work first with \((Y,q,d)\) as in the definition. Mixed masses of
\(T_v\) and \(q\) are all \(V_q=\int_Yq^d\), by the mass identity above.
Thus \(E_q(0)=0\) and \(E_q(v+b)=E_q(v)+bV_q\).

Let \(v_t=(1-t)w+tv\) for bounded \(q\)-psh endpoints, and write
\(h=v-w\). Positivity and boundedness hold for \(0\leq t\leq1\).
Multilinearity makes each term defining \(E_q(v_t)\) a finite polynomial
in \(t\) with integrable bounded weights; differentiation is therefore
algebraic, not a differentiation of an uncontrolled weak limit. For the
term with \(j\geq1\), apply (IBP) with \(a=v_t\), \(b=h\), and
\(\Theta=T_{v_t}^{j-1}q^{d-j}\). These satisfy every antecedent listed
above. It replaces \[\int_Y v_t\,dd^ch\,T_{v_t}^{j-1}q^{d-j}
\quad\hbox{by}\quad
\int_Y h\,(T_{v_t}-q)T_{v_t}^{j-1}q^{d-j}.\] Consequently, on the open
segment, with one-sided derivatives at endpoints, \[\frac{d}{dt}E_q(v_t)
=\frac1{d+1}\sum_{j=0}^d\left[(j+1)\int_YhT_{v_t}^jq^{d-j}
-j\int_YhT_{v_t}^{j-1}q^{d-j+1}\right]
=\int_YhT_{v_t}^d. \tag{9}\] The second term at \(j=0\) is absent. The
finite sum telescopes.

Integrating (9) and expanding \(T_{v_t}^d\) gives \[E_q(v)-E_q(w)
=\frac1{d+1}\sum_{j=0}^d\int_Y(v-w)T_v^jT_w^{d-j}
=R_q(v,w), \tag{10}\] because
\(\int_0^1\binom dj t^j(1-t)^{d-j}\,dt=1/(d+1)\); integration by parts
in the ordinary one-variable integral, or the elementary beta-integral
formula, proves this coefficient. Equation (10) proves antisymmetry and
the exact cocycle \[R_q(v,w)+R_q(w,z)=R_q(v,z).\] For ordered endpoints
\(w\leq v\), (9) is nonnegative, proving monotonicity. If
\(\delta=\|v-w\|_\infty\), then \(w-\delta\leq v\leq w+\delta\).
Monotonicity and the shift law give
\[|E_q(v)-E_q(w)|\leq\delta V_q. \tag{11}\] This holds for arbitrary
bounded, possibly discontinuous, endpoints.

Differentiating the polynomial in (9) once more and applying (IBP) with
\(a=b=h\) and \(\Theta=T_{v_t}^{d-1}\) gives \[\frac{d^2}{dt^2}E_q(v_t)
=d\int_Yh\,dd^ch\,T_{v_t}^{d-1}
=-d\int_Ydh\wedge d^ch\wedge T_{v_t}^{d-1}\leq0. \tag{12}\] The gradient
square is a finite positive measure by the cited bounded theorem, so
(12) proves concavity throughout the bounded convex domain. It also
specifies exactly what the second variation means for nonsmooth bounded
endpoints.

There is an additional useful first-variation scope. If \(t\mapsto v_t\)
is a \(C^1\) map into the Banach space \(C(Y)\) with the sup norm, and
every \(v_t\) is \(q\)-psh, then
\[\frac{d}{dt}E_q(v_t)=\int_Y\dot v_t T_{v_t}^d. \tag{13}\] Here
\(\dot v_t\) is continuous. To verify it, first note that uniform
convergence of bounded \(q\)-psh \(a_l\) to \(a\) implies weak
convergence of each mixed measure \(T_{a_l}^jT_a^{d-j}\) to \(T_a^d\).
For a smooth test \(f\), telescope one differing factor at a time and
use BT recursion and Stokes to bound the difference by
\[j\|a_l-a\|_\infty C_f V_q,\] where \(-C_fq\leq dd^cf\leq C_fq\). All
mixed masses are \(V_q\). Smooth test approximation and positivity
extend this convergence to continuous tests. Apply (10) to
\(v_{t+h},v_t\), divide by \(h\), use the uniform convergence
\((v_{t+h}-v_t)/h\to\dot v_t\) and the just proved mixed weak
convergence, and obtain (13). No conclusion with an arbitrary
discontinuous velocity or a merely measurable path is inferred. For
smooth paths of smooth admissible potentials this specializes to the
usual formula; only affine bounded paths have the second-variation
formula (12) asserted here.

All these arguments apply to \(q_\epsilon\) on \(X\) for each fixed
positive \(\epsilon\), and to \(\beta\) on \(B\). They give the exact
reference, shift, cocycle, monotonicity, concavity, Lipschitz, and
variational assertions of the theorem.

\subsection{Leading energy and arbitrary uniform perturbative
approximants}\label{leading-energy-and-arbitrary-uniform-perturbative-approximants}

Put \(H=\theta+dd^cU\), a positive bounded-potential current by
holomorphic pullback. Products of more than \(\nu\) factors chosen from
\(H\) and \(\theta\) vanish by the local pullback argument. Expand
\[E_{q_\epsilon}(U)
=\frac1{n+1}\sum_{j=0}^n\int_XU\,(H+\epsilon\omega_\epsilon)^j
(\theta+\epsilon\omega_\epsilon)^{n-j}.\] Every surviving term with
\(a\) factors \(H\) and \(b\) factors \(\theta\) has \(a+b\leq\nu\) and
power \(\epsilon^{n-a-b}\). The coefficient at the least possible power
\(\epsilon^m\) is \[\frac1{n+1}\sum_{a=0}^\nu
\left[\sum_{j=0}^n\binom ja\binom{n-j}{\nu-a}\right]
\int_XU H^a\theta^{\nu-a}\omega_\epsilon^m.\] The bracket equals
\(\binom{n+1}{\nu+1}\). To prove it, count subsets of size \(\nu+1\) of
\(\{1,\ldots,n+1\}\) by their \((a+1)\)-st element \(j+1\); the earlier
and later choices are precisely the two binomial factors. A binomial
factor outside its admissible range means zero. Since
\[\frac1{n+1}\binom{n+1}{\nu+1}=\frac K{\nu+1},\] the exact weighted
projection (7) gives leading term \(s_\epsilon E_\beta(\varphi)\).

All remaining terms have order at least \(\epsilon^{m+1}\). Their
absolute integrals are bounded by
\[\|\varphi\|_\infty\int_XH^a\theta^b\omega_\epsilon^{n-a-b}
=\|\varphi\|_\infty\int_X\theta^{a+b}\omega_\epsilon^{n-a-b}.\] The
equality follows from the cohomological mass identity, including when
\(a+b=0\). There are finitely many terms, and the smooth auxiliary limit
bounds their coefficient integrals uniformly. Thus
\[E_{q_\epsilon}(U)=s_\epsilon E_\beta(\varphi)
+O(\|\varphi\|_\infty\epsilon^{m+1}). \tag{14}\] This is a
bounded-potential calculation using weighted projection on all fibres,
not the earlier smooth fixed-class formula.

Now apply (11), with \(q=q_\epsilon\), to \(u_\epsilon\) and \(U\):
\[|E_{q_\epsilon}(u_\epsilon)-E_{q_\epsilon}(U)|
\leq\delta_\epsilon\int_Xq_\epsilon^n. \tag{15}\] By (8),
\(\int_Xq_\epsilon^n/s_\epsilon=V_B+O(\epsilon)\). Dividing (14)--(15)
by \(s_\epsilon\) gives the quantitative conclusion
\[\left|\frac{E_{q_\epsilon}(u_\epsilon)}{s_\epsilon}-E_\beta(\varphi)\right|
\leq\delta_\epsilon(V_B+O(\epsilon))
+O(\|\varphi\|_\infty\epsilon). \tag{16}\] The constants in the \(O\)
terms are bounded for the stipulated auxiliary family. This proves (1)
with no epsilon-relative convergence rate. Apply it twice and use (10)
to prove (2). One may use any bounded reference \(\rho\) on the base and
approximating reference family: the limit then is
\(R_\beta(\varphi,\rho)\) rather than \(E_\beta(\varphi)\).

The hypotheses admit genuinely non-descending families and convergence
slower than \(\epsilon\). For example, with \(\varphi=0\), choose a
nonconstant smooth base function \(f\) and a sufficiently small smooth
\(\zeta\) on \(X\) that varies along a regular fibre and satisfies
\(\omega+dd^c\zeta>0\). Then
\[u_\epsilon=\sqrt\epsilon\,(f\circ\pi)+\epsilon\zeta\] is
\(\theta+\epsilon\omega\)-psh for all sufficiently small \(\epsilon\),
because \(\beta+\sqrt\epsilon\,dd^cf>0\), and is not a pullback. Its
uniform error has order \(\sqrt\epsilon\) (choose \(f\) with nonzero
norm). Existence of \(\zeta\) follows by taking a small smooth bump
varying on the positive-dimensional smooth regular fibre. This example
is only an admissibility illustration; the proof does not limit
approximants to this form.

\subsection{Supporting measure estimate and its role in the
energy}\label{supporting-measure-estimate-and-its-role-in-the-energy}

Write \(A=q_\epsilon+dd^cu_\epsilon\),
\(D=q_\epsilon+dd^cU=H+\epsilon\omega_\epsilon\). Multilinearity in the
bounded BT sense yields \[A^n-D^n=dd^c(u_\epsilon-U)\wedge
\sum_{j=0}^{n-1}A^jD^{n-1-j}. \tag{17}\] For a real smooth test \(f\) on
\(B\), put \(F=f\circ\pi\) and choose \(C_f\) such that
\(-C_f\beta\leq dd^cf\leq C_f\beta\). Pullback gives the same inequality
with \(\theta\) on \(X\), even where it is degenerate. For each positive
mixed current \(S_j=A^jD^{n-1-j}\), apply (IBP) with \(a=u_\epsilon-U\),
\(b=F\), \(\Theta=S_j\); all functions and currents meet its bounded
antecedents. Equivalently for this smooth test one may use the recursive
current identity and Stokes. It follows that
\[\left|\int_X F(A^n-D^n)\right|
\leq n C_f\delta_\epsilon I_\epsilon,\qquad
I_\epsilon=\int_X\theta q_\epsilon^{n-1}. \tag{18}\] Indeed each
\(\int_X\theta S_j\) equals \(I_\epsilon\) by the cohomological mass
identity. Its expansion starts at the required order:
\[I_\epsilon=\sum_{k=0}^{\nu-1}\binom{n-1}k\epsilon^{n-1-k}
\int_X\theta^{k+1}\omega_\epsilon^{n-1-k}
=\binom{n-1}{\nu-1}c_\epsilon V_B\epsilon^m
+O(\epsilon^{m+1}). \tag{19}\] Thus
\(I_\epsilon/s_\epsilon=(\nu/n)V_B+O(\epsilon)\), and the scaled error
in (18) is \(O_f(\delta_\epsilon)\), not
\(O_f(\delta_\epsilon/\epsilon^m)\).

The exact pullback expansion is \[\pi_*D^n=s_\epsilon C_\varphi^\nu+
\sum_{k=0}^{\nu-1}\binom nk\epsilon^{n-k}
\pi_*(H^k\omega_\epsilon^{n-k}). \tag{20}\] The remainder is a positive
measure of total mass \(O(\epsilon^{m+1})\), by the cohomological mass
identity; (6) supplies the exact leading term across all critical
values. After dividing by \(s_\epsilon\), its total variation is
\(O(\epsilon)\). Combining (18)--(20) proves convergence on smooth
tests. The scaled positive measures have total masses
\[\frac{\int_X A^n}{s_\epsilon}
=\frac{\int_Xq_\epsilon^n}{s_\epsilon}=V_B+O(\epsilon),\] so uniform
smooth approximation of any continuous test gives the asserted weak
convergence (3). Total-variation convergence is proved only for the
exact pullback family \(D\), not for arbitrary \(A\).

This explains the compatibility with the first variation of the base
energy: (9) and (13) pair the raw measure \(C_\varphi^\nu\) with the
allowed base velocities. There is no passage to a derivative of an
arbitrary perturbative family \(u_\epsilon(t)\) without additional
time-regularity hypotheses. Normalizing the positive-\(\epsilon\) mass
to one would instead yield \(C_\varphi^\nu/V_B\); it is not the
normalization in this theorem and supplies no reconstruction of \(V_B\)
from I-equivalence.

\subsection{Changes of references and auxiliary
directions}\label{changes-of-references-and-auxiliary-directions}

For two auxiliary Kähler forms in any classes, apply (1)--(3) with their
respective constants \(c_\omega\) from (5). Both energy limits are
exactly \(E_\beta(\varphi)\), and both relative limits exactly
\(R_\beta(\varphi,\psi)\). Thus the auxiliary class as well as its
representative disappears after the explicit transverse calibration.
Formula (5) separately explains why \(c\) changes with the class. This
is independence within the fixed supplied fibration and base reference,
not independence of the choice of carrier \(\pi\).

For smoothly varying \(\omega_\epsilon\) with Kähler limit, every
displayed remainder is uniform, and \(c_\epsilon\to c_{\omega_0}>0\).
Hence both instantaneous and limiting-\(c\) calibrations give the same
result. Any reparameterization of \(\epsilon\downarrow0\) inherits the
statement. No claim is made for auxiliary forms degenerating out of the
Kähler cone, unbounded coefficient integrals, or unrelated perturbations
that do not have the stated form.

Next let \(\beta'=\beta+dd^ch\) be a cohomologous Kähler representative,
with smooth real \(h\); on a compact Kähler manifold such a potential
exists by the \(dd^c\) lemma, or the displayed relation may simply be
included as the specified change. The \(dd^c\) lemma in the needed form
says that a real smooth closed \((1,1)\) form with zero de Rham
cohomology class is \(dd^c\) of a real smooth function on a compact
Kähler manifold; citation: P. Deligne, P. Griffiths, J. Morgan, D.
Sullivan, \emph{Real homotopy theory of Kähler manifolds}, Invent. Math.
29 (1975), 245--274, the \(dd^c\) lemma. It is used only for
interpreting arbitrary cohomologous representative changes; the formulas
themselves assume the exhibited \(h\).

Set \(\varphi'=\varphi-h\) and \(\psi'=\psi-h\). Then
\(C_{\varphi'}'=C_\varphi\), \(C_{\psi'}'=C_\psi\), and
\(\varphi'-\psi'=\varphi-\psi\). The exact relative formula gives
\[R_{\beta'}(\varphi-h,\psi-h)=R_\beta(\varphi,\psi). \tag{21}\] Since
\(h\) is bounded \(\beta\)-psh, (10) yields the reference correction
\[E_{\beta'}(\varphi-h)=R_\beta(\varphi,h)
=E_\beta(\varphi)-E_\beta(h). \tag{22}\] Also \(\int_B(\beta')^\nu=V_B\)
by Stokes. Adding a constant to \(h\) changes the associated potential
and reference by the corresponding shift in (4); it does not give an
absolute reference-free energy.

These identities already hold at positive \(\epsilon\). Write
\(H_h=h\circ\pi\), \(q'_\epsilon=q_\epsilon+dd^cH_h\), and
\(u'_\epsilon=u_\epsilon-H_h\). Both backgrounds are Kähler, the total
current of \(u'_\epsilon\) is that of \(u_\epsilon\), and the bounded
relative formula shows \[E_{q'_\epsilon}(u'_\epsilon)
=E_{q_\epsilon}(u_\epsilon)-E_{q_\epsilon}(H_h).\] Its calibrated limit
is precisely (22). If a common physical reference current is retained,
transform its potential too, and (21) gives exact invariance of the
relative energy. This states the reference/cocycle scope explicitly
rather than asserting representative independence of an absolute energy
with different zero references.

\subsection{Optional continuous descent with connected
fibres}\label{optional-continuous-descent-with-connected-fibres}

This paragraph is not required for (1)--(3), which begin with a supplied
bounded base potential. Assume additionally that every fibre is
connected and let \(W\) be a continuous \(\theta\)-psh function on
\(X\). Then there exists a unique continuous \(\beta\)-psh \(\varphi\)
on \(B\) with \(W=\varphi\circ\pi\).

Here the sole additional geometric foundation is complex analytic
resolution: a compact irreducible reduced complex analytic space admits
a proper surjective holomorphic map from a smooth connected compact
complex manifold, which is an isomorphism over a dense nonsingular open
subset. This is the analytic desingularization consequence of E.
Bierstone and P. D. Milman, \emph{Canonical desingularization in
characteristic zero by blowing up the maximum strata of a local
invariant}, Invent. Math. 128 (1997), 207--302, Main Theorem and its
complex analytic category. We invoke it only for irreducible components
of the reduced compact fibre, not for resolving the map \(\pi\).
Compactness of the resolution follows from properness; connectedness
follows by taking the resolving component dominating the irreducible
space.

Fix a fibre and an irreducible component \(Z\) of its reduction. In
local base coordinates \(\beta=dd^c\rho\), so \(W+\rho\circ\pi\) is psh
near the fibre. On a resolution of \(Z\), its pullback differs from
\(W\) by the constant \(\rho(b)\), hence \(W\) pulls back to an ordinary
psh function on a connected compact complex manifold. Holomorphic
pullback is legitimate because \(W\) is continuous and finite. The
ordinary maximum principle, applied at a maximum on this smooth compact
manifold, makes it constant. Surjectivity of the resolution makes \(W\)
constant on \(Z\). There are finitely many fibre components; components
joined by intersections have the same value, and connectedness of the
reduced fibre makes their intersection graph connected. Thus \(W\) is
constant on the entire fibre, including singular and higher-dimensional
components.

A continuous surjection from a compact space to a Hausdorff space is a
quotient map. Therefore the uniquely defined fibrewise function
\(\varphi\) is continuous. It remains to prove psh positivity on the
base, without assuming a submersion at critical values. For any fixed
auxiliary \(\omega\), bounded multiplication of order-zero currents and
Stokes give \[\pi_*\big((\theta+dd^cW)\wedge\omega^m\big)
=c_\omega\beta+dd^c\pi_*(W\omega^m)
=c_\omega(\beta+dd^c\varphi).\] The left side is a positive current, and
\(c_\omega>0\), so \(\beta+dd^c\varphi\geq0\) distributionally. A
continuous function whose sum with a smooth local potential has positive
\(dd^c\) is psh: on smaller coordinate balls convolve it with radial
positive approximate identities, obtaining smooth psh functions
converging locally uniformly, and apply the submean inequalities on
complex lines in the limit. This proves the assertion through all
critical values. No descent statement for arbitrary unbounded potentials
or disconnected fibres is made.

\subsection{What has and has not been
resolved}\label{what-has-and-has-not-been-resolved}

The result extends the supplied old comparison result in the precise
requested directions: arbitrary supplied maps with singular fibres;
bounded and possibly discontinuous base potentials; non-descending
bounded perturbative potentials; uniform convergence with no
\(o(\epsilon)\) requirement; complete removal of the auxiliary Kähler
class after its own transverse calibration; smoothly converging
auxiliary directions; relative reference corrections; and rigorously
specified bounded variational scopes. It uses the raw positive base
mass, never the zero top mass of \(\alpha\).

The additional structure is substantial. A general pseudoeffective
zero-volume class need not be the pullback of a Kähler class under any
supplied holomorphic map, and need not admit a smooth semipositive
representative. No canonical fibration or minimal reference scale for
that general class has been constructed. Independence under changing the
supplied map/carrier is not proved. No unbounded prescribed P-type, E1
completion, or arbitrary degenerate auxiliary path is included. All
bounded potentials have the minimal bounded singularity type relative to
these smooth semipositive backgrounds (the smooth potential \(0\) is
already admissible, and every quasi-psh potential is bounded above on a
compact manifold); this is not a theorem retaining arbitrary unbounded
P-data. The optional descent statement covers continuous potentials
only. The original general P15.1 remains open.

The catalogue's I-equivalence obstruction, its published class-level
convex-body endpoint, and its repaired unpublished mixed-volume results
are not used. Their descriptions are context, not independently
established premises. Likewise earlier model verdicts in the comparison
papers do not certify this result.

\section*{Review and supporting evidence}

The independent Xean native Codex review used \texttt{gpt-6.1-sol} with
\texttt{xhigh} reasoning. The entire report, 6 premises and
9 primary-source passages were read and bound to the unchanged
author argument, TeX source and partial status. The reviewer reconstructed
the new applications of the standard bounded Bedford--Taylor calculus,
integration by parts, cohomological mass identities and geometric
foundations. No substantive mathematical repair was required.

The complete author argument and both actual canonical native notes were
read. Six supplied primary bodies were checked against their exact LF
bytes and source hashes. Both diagnostic and two-pass document builds
were repeated in fresh copies, producing TeX byte-identical to the frozen
author source. Publication edits update review provenance, citations and
explicit math-alphabet grouping. The supporting packet preserves the
full frozen argument, review, canonical notes, source provenance and read
records; native editor compilation and PDF exports are recorded separately.

This review establishes the stated special-case result. Same-model review,
document compilation and file integrity are not formal verification or
priority certification. The original general problem P15.1 remains open;
no general carrier construction, minimum-reference theorem, unbounded
prescribed-type extension or arbitrary degenerate perturbation path is
claimed.

\end{document}
